\section{Introduction}
\label{introduction}

\subsection{Background and Motivation}
\label{sec:intro-background}

Requirements are the primary interface between stakeholder intent and downstream software engineering activities.
Recent work has emphasized the central role of requirements in AI-driven software engineering~\cite{robinsonRequirementsAreAll2025}.
However, natural-language requirements remain difficult to analyze because their structure is implicit, their semantics are mixed, and their domain assumptions are often left unstated.

\subsection{Challenges of Requirement Structure Modeling}
\label{sec:intro-challenges}

This paper focuses on requirement structure modeling rather than on later development activities.
We identify four core challenges:
\begin{enumerate}[leftmargin=*]
  \item Natural-language requirements mix functions, constraints, roles, data elements, and exception conditions in the same text.
  \item Requirement semantics are ambiguous when stakeholder vocabulary, system behavior, and operational conditions are not explicitly separated.
  \item Existing LLM-based requirement analysis often produces fluent text but lacks a stable structured representation.
  \item Requirement representations must remain inspectable and refinable by humans to support practical requirements engineering.
\end{enumerate}

\subsection{Multi-Agent Requirement Parsing}
\label{sec:intro-multi-agent}

We study a lightweight multi-agent parsing workflow in which different agents focus on elicitation, semantic analysis, and structure construction.
The goal is not to build a large agent platform, but to examine whether role-separated collaboration can improve the transformation from informal requirements to a structured requirement graph.

\subsection{Contributions}
\label{sec:intro-contributions}

The intended contributions of this paper are:
\begin{enumerate}[leftmargin=*]
  \item A research formulation of knowledge-guided requirement structure modeling for natural-language requirements.
  \item A multi-agent workflow for requirement semantic decomposition and structured representation construction.
  \item A requirement graph representation that organizes functional requirements, constraints, roles, data elements, and exception information.
  \item An evaluation design focused on structure completeness, semantic consistency, graph quality, and representation accuracy.
\end{enumerate}
